% !TeX root = ../main.tex
% Operation and asset management.
\chapter{Operation and Asset Management}
\label{ch:assets}
\label{ch:delivery-assets}

An operator inherits a structure whose behaviour depends on ground conditions recorded, if at all, in construction documents nobody has read since handover. The asset must then be inspected, assessed, repaired and eventually repurposed over a century, and how well that can be done is largely decided by what was written down at the end of construction.

\section{Condition, inspection and assessment}
\label{sec:condition-inspection}

\subsection{Transition to operation and deterioration}
\label{sec:transition-deterioration}

Handover is where most asset management fails, and it fails through omission
rather than error. The operator receives drawings and an operating manual, and
frequently does not receive the ground model, the as-built ground conditions
chainage by chainage, the residual risk register, the monitoring baseline, or
the design assumptions on which durability depends. Those documents exist ---
they were produced during design and construction --- and they are archived
rather than transferred.

The consequence appears decades later, when an assessment must be made and
nobody can establish what the structure was designed for. A handover package
should therefore include the ground model
(Chapter~\ref{ch:ground-modelling}), the as-built record against the baseline,
the residual risks with their intended management, the monitoring baseline,
and the maintenance assumptions the durability design relies on.

Deterioration proceeds by a small number of mechanisms. Water ingress is the
common driver: it carries the chemistry that attacks concrete, delivers
chlorides to reinforcement, leaches calcium from the matrix, and in cold
climates drives freeze--thaw damage. Reinforcement corrosion follows
depassivation by chloride or carbonation and is the dominant structural
deterioration mode. Chemical attack from aggressive groundwater continues at a
rate set by the exposure assessed in Chapter~\ref{ch:hydrogeology}. Slow ground
movement --- creep, adjacent construction, groundwater change --- loads the
lining in ways not anticipated. And in older masonry and cast-iron tunnels,
mortar loss and joint deterioration progress steadily.

\subsection{Inspection and condition assessment}
\label{sec:inspection-condition}

An inspection regime combines frequent general inspections, largely visual and
carried out during available access, with less frequent principal inspections
that are detailed, close-up and instrumented. Intervals are set by consequence
and by observed deterioration rate rather than by convention.

Access is the binding constraint. A rail tunnel can be inspected only during
possessions; a road tunnel requires closures; a pressurised water tunnel must
be dewatered, which may be possible only every several years and at
substantial cost in lost service. Inspection technology has advanced largely
in response to this: mobile laser scanning and imaging systems that survey a
whole tunnel at line speed convert a multi-shift manual inspection into a
single pass.

Beyond visual survey, ground penetrating radar locates voids behind the lining
and measures its thickness; impact echo and ultrasonic methods detect
delamination; thermography reveals moisture and voids; and automated defect
detection from imagery classifies cracking and spalling with a consistency
manual inspection does not achieve.

Condition rating converts observation into a prioritisation. What matters is
that the rating scale distinguishes defects that threaten structural capacity
from those that threaten durability from those that are merely unsightly,
because the three warrant entirely different responses and a single composite
score conceals the distinction.

Predictive maintenance is the direction of travel. Surface observation,
internal inspection and physical modelling combined can forecast condition
rather than record it, which moves intervention ahead of loss of function.
This depends on inspection data returning to the asset model of
Chapter~\ref{ch:digital-twin} rather than accumulating in inspection reports,
and it is the operational half of the lifecycle argument set out in
Chapter~\ref{ch:challenges}.

\subsection{Assessment of structures in service}
\label{sec:assessment-in-service}

Assessing an existing tunnel is harder than designing a new one, because the
geometry, the material properties and the loads are all uncertain and the
structure cannot be taken apart to find out. Masonry, cast iron and early
concrete linings were built without the records a modern project produces, and
their thickness may vary along the tunnel in ways no drawing shows.

The assessment therefore begins with investigation rather than analysis:
establishing thickness and construction by radar and coring, material
properties by testing recovered samples, and the ground behind the lining ---
whether it is in contact at all --- by radar and probe drilling. Voids behind
a lining are common in older tunnels and change the load path fundamentally.

Back-analysis from observed condition is often the most informative step. A
tunnel that has stood for a century with no distress has demonstrated
something about its loading that no calculation will improve on; a tunnel with
a characteristic crack pattern has revealed the mechanism acting on it.
Assessing such a structure against modern code loading, without that evidence,
frequently condemns a tunnel that is performing perfectly well.

Long-term instrumentation supports this, particularly where new construction
passes nearby and the existing structure must be shown to be unaffected
(Chapter~\ref{ch:settlement-damage}).

\section{Intervention and whole-life management}
\label{sec:intervention-whole-life}

\subsection{Repair, strengthening and upgrading}
\label{sec:repair-strengthening}

Repair should address the cause rather than the symptom, and for tunnel
linings the cause is usually water. Injection grouting seals cracks and voids;
patch repair reinstates spalled concrete; drainage rehabilitation restores a
system that has clogged; cathodic protection arrests reinforcement corrosion
where chloride contamination is widespread. Sealing a leak without restoring
the drainage that should have carried the water away simply relocates it.

Strengthening becomes necessary where capacity is genuinely deficient. Sprayed
concrete relining is the common solution, adding a structural layer bonded to
the existing lining; segmental or sleeve relining inserts a new structure
inside the old. Both consume clearance, and in an operating tunnel the
clearance envelope is usually the constraint that decides the method --- which
is why the allowance for future relining belongs in the original internal
diameter (Chapter~\ref{ch:linings}).

Upgrading legacy tunnels to modern standards is a distinct problem. Nineteenth
and twentieth century tunnels were built to clearances, fire standards and
egress provisions that no longer suffice, and bringing them into compliance
within an immovable bore, often with heritage constraints, is among the more
difficult problems in the field. The honest options are frequently to accept a
lower service standard with compensating measures, or to build a new tunnel.

Restoration should be evaluated as a design decision, on functionality
recovered, time, cost, constructability, risk and carbon together --- not on
repair cost alone. A cheaper repair requiring three months more closure may be
the more expensive option once lost service is counted, which is the resilience
argument of Chapter~\ref{ch:sustainability} applied to a specific decision.

\subsection{Operational safety and whole-life management}
\label{sec:operational-whole-life}

Operational safety systems --- ventilation and smoke control, fire detection
and suppression, emergency egress and cross-passage spacing, communications
and incident management --- are designed with the tunnel and constrain it. Fire
life safety in particular determines cross-passage spacing, ventilation
capacity and, through them, the excavated volume. It is not an add-on to a
completed geometry.

Whole-life cost management requires a deterioration model, a set of
intervention options with costs and effects, and an optimisation over the
remaining life. Frameworks such as the ISO 55000 series provide the structure.

The persistent difficulty is not technical. Maintenance is funded from
operating budgets and capital works from capital budgets, and the two are
rarely traded against each other, so deferred maintenance accumulates until it
becomes a capital problem. An asset whose durability design assumed drains
flushed every five years, in an organisation that funds flushing every
fifteen, has been designed against a maintenance regime that does not exist.
Recording the maintenance assumption explicitly at handover is the designer's
only defence.

\subsection{Decommissioning and repurposing}
\label{sec:decommissioning}

Underground assets are rarely demolished. The options at end of service are
abandonment with the openings sealed, infilling to prevent collapse, monitoring
in perpetuity, or repurposing.

Repurposing is increasingly common and increasingly valuable: disused rail
tunnels become cycleways and utility routes, mine workings become pumped
storage and data centres, and urban tunnels are reused for services. The
engineering question is whether the structure can be assessed and brought to a
condition suitable for its new duty, which returns to the assessment problem
above --- and to the handover documentation that determines whether the
assessment is feasible at all.

Abandonment is not costless. A sealed tunnel remains a void in the ground with
the subsidence potential described in
Chapter~\ref{ch:settlement-damage}, and the obligation to record its position
and condition persists indefinitely.

\section{Digital continuity and security}
\label{sec:digital-continuity}

An underground asset is designed for a century and its information for
whatever the current software supports. That mismatch is not a minor
administrative problem: an operator who cannot open the ground model, cannot
locate the as-built record and cannot trust the monitoring feed is operating
the asset blind, however sophisticated the delivery was. Two threats produce
that outcome. One is deliberate, and arrives through the connectivity that
makes a live twin possible. The other is entirely passive and simply requires
that nobody does anything for fifteen years.

\subsection{The operational record}
\label{sec:operational-record}

The inspection regime of Section~\ref{sec:inspection-condition} generates a
record whose value is almost entirely comparative. A crack width means little
in isolation and a great deal alongside the same crack measured five years
earlier; a damp patch matters when it was dry at the last inspection. The
record is therefore an instrument, and like any instrument it produces nothing
useful if its datum moves.

Three requirements follow. Observations must be located against a stable
reference, which is what the as-built survey of
Chapter~\ref{ch:movement-monitoring} provides --- a chainage and a position on
a surveyed surface, rather than a description that depends on who wrote it.
They must be recorded to a consistent scheme, so that a change in
classification between inspection contracts does not appear as a change in
condition. And the raw observation must be retained alongside the assessment,
because assessment criteria are revised and an old rating cannot be
reinterpreted while the measurement behind it can.

The handover record described in Chapter~\ref{ch:delivery} is the origin of
this series, and its quality determines what the operator can ever conclude.
An asset handed over with a dimensioned as-built model begins its
deterioration record with a known starting condition; one handed over with a
set of design drawings begins with an assumption.

\subsection{Security of connected systems}
\label{sec:cyber-security}

A twin that updates from live instrumentation is a connected control-adjacent
system attached to critical infrastructure, and it inherits the exposure that
implies. The realistic concern is not the theatrical one. Remote seizure of
ventilation or fire systems is the scenario that gets discussed; corruption of
the monitoring record is the scenario that is both easier and more damaging,
because a monitoring system reporting normal readings while the ground moves
defeats the observational method precisely when it is needed, and does so
without anyone noticing that anything has happened.

The established engineering response applies here without modification.
Operational technology should be separated from corporate networks, with
traffic between them constrained to what is necessary and inspected.
Monitoring data should be authenticated at the source, so that a reading can
be tied to the instrument that produced it. Remote access should be
time-limited and individually attributable, which is chiefly a discipline
problem: shared vendor credentials for maintenance access are the most common
route in. And the trigger levels of
Section~\ref{sec:triggers-interpretation} should be verifiable by
a manual reading, so that automated monitoring can be checked against
something that cannot be altered remotely.

Where an operator's response depends on the system, the failure of the system
must itself have a response. A documented degraded mode --- what is measured
manually, at what frequency, and who decides --- belongs in the operating
regime alongside the instrumentation, and it is also the answer to the
commoner and duller failure in which the data link simply stops working.

\subsection{Data longevity}
\label{sec:data-longevity}

The passive threat is obsolescence, and it is more likely to destroy the value
of a project's information than any attack. A tunnel commissioned today may be
inspected in 2080. The model will by then have outlived several generations of
the software that made it, the storage it sits on, the organisation that
delivered it and everyone who understood it.

Four measures address this, and none is technically difficult. Store the
information in open, documented formats, accepting the loss of proprietary
features as the price of being readable. Keep the underlying data as well as
the model, since data in a plain format can be rebuilt into whatever tool
exists later while a model cannot be reverse-engineered into data. Record the
context --- what the units are, what the coordinate system is, what the
parameters mean, what was assumed --- because an undocumented dataset is
unusable long before it is unreadable. And migrate deliberately at defined
intervals rather than trusting that the format will still open, verifying that
the migrated copy is complete rather than assuming it.

The obligation is the operator's, but the possibility is created or destroyed
at handover. A specification that requires open formats and accompanying data
costs little at the time of writing and is the difference between an asset
whose history is available in fifty years and one whose history exists only in
the memory of people who have retired.
